% ===== §13 =====
\section{The Dialectic of the Material and the Ideal}
\label{sec:material}

\noindent
The antithesis closes with the dialectic on which the whole framework balances, and which the two preceding sections have each leaned on without stating: the relation between the material and the ideal in culture. The account of culture as the exteriorization of drive, and before it the inversion that made culture a generation rather than a deposit, both require that culture be at once produced by material conditions and irreducible to them---an ideal form with its own efficacy, yet never floating free of the material ground that conditions it. This double requirement is the site of the two vulgarizations the paper must refuse, and refusing both is what its position on the material and the ideal amounts to.

\subsection*{Two vulgarizations refused}

The first vulgarization is \emph{mechanical determinism}: the reduction of culture to a passive reflection of the economic base, a mere epiphenomenon that registers material conditions and does nothing back. On this view the culture a couple generates would be simply the shadow their material circumstances cast, and to explain the circumstances would be to explain the culture entirely. This is the error against which Engels warned in his later correspondence, and the warning is worth stating in his own terms because it is often forgotten by those who invoke historical materialism most loudly. The ultimately determining element in history, Engels wrote to Bloch in 1890, is the production and reproduction of real life; but to twist this into the claim that the economic element is the \emph{only} determining one is to transform the proposition into a meaningless, abstract, senseless phrase \citep{engels1890bloch}. The economic situation is the base, but the elements of the superstructure exert their own influence on the course of struggles and in many cases preponderate in determining their \emph{form}. Culture, that is, has a \emph{relative autonomy} and a genuine reciprocal action upon the conditions that produce it; it is determined in the last instance, not in the first, and the space between the last instance and the first is exactly the space in which culture does its own work. This relative autonomy is not a concession that weakens the materialist thesis; it is what gives the thesis a culture worth explaining rather than a shadow.

The second vulgarization is the opposite error, \emph{idealism}: the treatment of culture as a self-generating order of meaning or spirit that spins itself out of its own resources, owing nothing essential to the material conditions of its production. This is the error the Hegelian line of the paper, taken alone, would tempt one toward---for if culture is spirit exteriorizing and knowing itself, one might forget that the exteriorization is always into a material medium and under material constraint, and imagine a self-movement of pure relational spirit generating culture out of itself. The framework refuses this as firmly as it refuses the first error. The dialectic of the registers has a material motor, as the meta-theoretical statement insisted: which signifiers are available to capture the Real, which fantasies are livable, which forms the drive can take, are set by the forces and relations of production, and no movement of the registers occurs except in and through a material medium that conditions it.

\subsection*{The dialectical position: constrained but not determined}

Between these two refusals stands the position the paper holds throughout, and it can now be stated as a thesis about the material and the ideal as such: culture is \emph{materially constrained but not materially determined}. The distinction is not a verbal hedge but the precise content of the dialectic, and it is the same distinction the intimacy series has maintained in its treatments of trust and of power, that a condition may bound a range of possibility without selecting a point within it. Material conditions set what is possible---they fix the outer limits of the forms culture can take, foreclose some exteriorizations of drive and enable others, supply the media and the relations within which any cultural generation must occur---and within those limits the generation is genuinely open, genuinely the work of the relational subjects whose drive and history it exteriorizes, and genuinely capable of reacting back upon the conditions that bounded it. The material is the floor and the walls; it is not the choreography. To say culture is constrained is to deny idealism: there is no cultural generation that escapes its material conditions, and the couple that imagines its private world a pure creation owing nothing to its circumstances is deceived. To say culture is not determined is to deny mechanism: the material conditions do not select which culture will be generated within the range they permit, and the couple's generation is a real generation and not the mere playing-out of a material script. Both denials are needed, and holding them together is the dialectic.

\subsection*{Why this closes the antithesis}

This dialectic closes the antithesis because it secures, at the level of the material and the ideal, exactly what the whole antithesis was for. The inversion made culture a generation; the account of drive gave the generation its content; and this dialectic gives the generation its correct relation to the material world, neither cut loose from it into idealism nor collapsed into it as its shadow. It is on this basis that the synthesis can define cultural emergence in a way that is at once materialist---bound to the forces and relations of production, historical rather than eternal---and non-reductive---a real generation with its own efficacy, carrying the exteriorized drive of relational subjects across the three registers. The synthesis can now be stated, for the ground it requires is complete: an inversion of the shared presupposition, a psychoanalytic account of what culture exteriorizes, and a dialectic of the material and the ideal that keeps the exteriorization honest. The definition of the emergence of culture in intimate relation follows.
